\begin{abstract}
Task offloading in edge-cloud Internet of Things (IoT) systems changes with task urgency, channel quality, edge queue pressure, cloud load, and fault state. Policies trained on mostly independent task samples can miss the behavior that appears during bursty arrivals, outage windows, queue carryover, and infeasible tasks. This article proposes a two-stage scheduling method for channel-aware offloading and resource management. Stage 1 uses a zone-aware Federated Double Deep Q-Network (\fedddqn) to choose between edge and cloud execution across non-identical edge zones. Rejection is modeled as an admission or infeasibility outcome, not as a third action. Stage 2 assigns CPU, memory, and bandwidth to edge-selected tasks through a rejection-aware residual allocator. The method is evaluated on \dataset{}, a realistic synthetic edge-cloud IoT benchmark with 100,000 tasks, temporal workload bursts, congestion, outage, jitter, edge degradation, and cloud maintenance. On the untouched temporal test split, \fedddqn{} achieved \FedLatency{} ms average latency, \FedSLA\% SLA satisfaction, \FedRejection\% rejection, and \FedEdgeUse\% edge usage. Its latency was \FedVsDDQN\% lower than centralized DDQN and \FedVsFLDDPG\% lower than an FL-DDPG-style baseline; the latency-reference Oracle reached \OracleLatency{} ms. On the allocator-valid edge-selected subset, the hand-designed rejection-aware demand rule gave the lowest proxy-target MAE, while residual learning with risk projection and capacity normalization gave the best efficiency score, lowest under-allocation, and zero capacity violation. The results indicate that the proposed pipeline improves latency and QoS balance under the evaluated edge-cloud workload conditions, while the latency-reference Oracle gap leaves clear room for further optimization.
\end{abstract}

\begin{keywords}
Double Deep Q-Network, edge-cloud computing, federated reinforcement learning, Internet of Things, quality of service, resource allocation, task offloading.
\end{keywords}
